\documentclass[12pt,a4paper]{article}

% --- PAQUETES ACTIVOS ---
\usepackage[utf8]{inputenc}
\usepackage[T1]{fontenc}
\usepackage[spanish]{babel}
\usepackage{csquotes}
\usepackage{lmodern}
\usepackage{microtype}
\usepackage{geometry}
\usepackage{setspace}
\usepackage{fancyhdr}
\usepackage{xcolor}
\usepackage{graphicx}
\usepackage{float}
\usepackage{caption}
\usepackage{subcaption}
\usepackage{booktabs}
\usepackage{tabularx}
\usepackage{longtable}
\usepackage{multirow}
\usepackage{array}
\usepackage{enumitem}
% \usepackage{tcolorbox}  % DESACTIVADO en modo redacción — reactivar para versión final impresa
\usepackage{hyperref}
\usepackage[style=apa, backend=biber]{biblatex}
\usepackage{titlesec}

% --- BIBLIOGRAFÍA ---
\addbibresource{TFM-ML-referencias.bib}

% --- ENTORNOS ESPECIALES (versión ligera para compilación rápida) ---
% NOTA: tcolorbox comentado durante redacción activa. Reactivar para versión final de impresión.
\newenvironment{pendiente}
  {\par\vspace{4pt}\noindent\textcolor{gray}{\textbf{[PENDIENTE]}}\par\begingroup\color{gray}\small}
  {\endgroup\par\vspace{4pt}}

\newenvironment{claimS}
  {\par\vspace{4pt}\noindent\textcolor{green!60!black}{\textbf{[CLAIM SUPPORTED]}}\par\begingroup}
  {\endgroup\par\vspace{4pt}}

\newenvironment{claimC}
  {\par\vspace{4pt}\noindent\textcolor{red!70!black}{\textbf{[CLAIM CONTESTED]}}\par\begingroup}
  {\endgroup\par\vspace{4pt}}

\newenvironment{claimU}
  {\par\vspace{4pt}\noindent\textcolor{orange!80!black}{\textbf{[CLAIM UNDETERMINED]}}\par\begingroup}
  {\endgroup\par\vspace{4pt}}

\newenvironment{notametod}
  {\par\vspace{4pt}\noindent\textcolor{blue!70!black}{\textbf{[NOTA METODOLÓGICA]}}\par\begingroup}
  {\endgroup\par\vspace{4pt}}

\newcommand{\fuentependiente}[1]{\textcolor{red}{[FUENTE PENDIENTE: #1]}}

% --- CONFIGURACIÓN ---
\geometry{left=3cm,right=2.5cm,top=3cm,bottom=3cm}
\onehalfspacing

\title{TFM: Sistemas de Transporte Continuo}
\author{Mario Lourido Regueira}
\date{\today}

\begin{document}

\maketitle

% ============================================================
% PÁGINAS PRELIMINARES (maquetación en InDesign, no numeradas)
% - Portada institucional (según formato EUDI)
% - Resumen / Abstract
% - Índice general
% - Índice de figuras y tablas
% - Glosario de términos y acrónimos
% ============================================================


% ============================================================
% CAPÍTULO 1 — INTRODUCCIÓN
% ============================================================

\section{Introducción}

\subsection{Presentación del proyecto}
\begin{pendiente}
\end{pendiente}

\subsection{Motivación y justificación}
\begin{pendiente}
\end{pendiente}

\subsection{Objetivos}
\begin{pendiente}
\end{pendiente}

\subsection{Alcance y limitaciones}
\begin{pendiente}
\end{pendiente}

\subsection{Metodología}
\begin{pendiente}
\end{pendiente}

\subsection{Estructura de la memoria}
\begin{pendiente}
\end{pendiente}


% ============================================================
% CAPÍTULO 2 — MARCO CONTEXTUAL
% ============================================================

\section{Marco Contextual}

\subsection{La ciudad de 24 horas: evolución del concepto}
\begin{pendiente}
\end{pendiente}

\subsection{Análisis PESTEL del transporte nocturno en Europa}

\subsubsection{Político}
\begin{pendiente}
\end{pendiente}

\subsubsection{Económico}
\begin{pendiente}
\end{pendiente}

\subsubsection{Social}
\begin{pendiente}
\end{pendiente}

\subsubsection{Tecnológico}
\begin{pendiente}
\end{pendiente}

\subsubsection{Ecológico}
\begin{pendiente}
\end{pendiente}

\subsubsection{Legal}
\begin{pendiente}
\end{pendiente}

\subsection{Mapa de stakeholders}
\begin{pendiente}
\end{pendiente}


% ============================================================
% CAPÍTULO 3 — INVESTIGACIÓN DE USUARIO Y MERCADO
% ============================================================

\section{Investigación de Usuario y Mercado}

\subsection{Definición y segmentación de usuarios}
\begin{pendiente}
\end{pendiente}

\subsection{Necesidades, barreras y expectativas por segmento}
\begin{pendiente}
\end{pendiente}

\subsection{Benchmarking de sistemas existentes}

\subsubsection{Vertical A --- Transporte nocturno urbano}
\begin{pendiente}
\end{pendiente}

\subsubsection{Vertical B --- Trenes nocturnos interurbanos}
\begin{pendiente}
\end{pendiente}

\subsubsection{Vertical C --- Conducción autónoma nocturna}
\begin{pendiente}
\end{pendiente}

\subsection{Vigilancia tecnológica}
\begin{pendiente}
\end{pendiente}


% ============================================================
% CAPÍTULO 4 — ESTRATEGIAS DE DISEÑO
% ============================================================

\section{Estrategias de Diseño}

\subsection{Valores del sistema}
\begin{pendiente}
\end{pendiente}

\subsection{Requisitos de diseño (QFD)}
\begin{pendiente}
\end{pendiente}

\subsection{Priorización MoSCoW}
\begin{pendiente}
\end{pendiente}

\subsection{DAFO por vertical}
\begin{pendiente}
\end{pendiente}

\subsection{Decisión de scope}
\begin{pendiente}
\end{pendiente}


% ============================================================
% CAPÍTULO 5 — CONCEPTO DE SISTEMA
% ============================================================

\section{Concepto de Sistema}

\subsection{Descripción del concepto}
\begin{pendiente}
\end{pendiente}

\subsection{Modelo de negocio}
\begin{pendiente}
\end{pendiente}

\subsection{Gobernanza del sistema}
\begin{pendiente}
\end{pendiente}

\subsection{Journey map del usuario}
\begin{pendiente}
\end{pendiente}

\subsection{Interacción humano-sistema}
\begin{pendiente}
\end{pendiente}


% ============================================================
% CAPÍTULO 6 — EVALUACIÓN DEL CONCEPTO
% ============================================================

\section{Evaluación del Concepto}

\subsection{Pliego de condiciones del sistema}
\begin{pendiente}
\end{pendiente}

\subsection{Matriz de afirmaciones-evidencia}
\begin{pendiente}
\end{pendiente}

\subsection{Valor diferencial de la propuesta}
\begin{pendiente}
\end{pendiente}

\subsection{Viabilidad y riesgos}
\begin{pendiente}
\end{pendiente}


% ============================================================
% CAPÍTULO 7 — CONCLUSIONES
% ============================================================

\section{Conclusiones}

\subsection{Síntesis de hallazgos}
\begin{pendiente}
\end{pendiente}

\subsection{Limitaciones del trabajo}
\begin{pendiente}
\end{pendiente}

\subsection{Líneas futuras}
\begin{pendiente}
\end{pendiente}


% ============================================================
% CAPÍTULO 8 — PRESUPUESTO DEL PROYECTO
% ============================================================

\section{Presupuesto del Proyecto}
\begin{pendiente}
\end{pendiente}


% ============================================================
% BIBLIOGRAFÍA
% ============================================================

\printbibliography


% ============================================================
% ANEXOS (no numerados)
% ============================================================

\section*{Anexo A: Fichas completas de benchmarking}
\begin{pendiente}
\end{pendiente}

\section*{Anexo B: Datos brutos de fuentes primarias}
\begin{pendiente}
\end{pendiente}

\section*{Anexo C: Glosario extendido}
\begin{pendiente}
\end{pendiente}

\section*{Anexo D: Registro de horas del proyecto}
\begin{pendiente}
\end{pendiente}


\end{document}
